%%%%%%%%%%  scan idx 212  |  folio 199  %%%%%%%%%%
\SGAPrintPage{199}%
\sgaStableTarget{sga6.e02.appendix-ii-definition-of-the-analytic-index-of}\section*{Appendix II\\ Definition of the analytic index of a relative elliptic complex}

In this appendix, we propose to show that the notion of a perfect complex leads to a simple and natural definition of the \textquotedblleft analytic index\textquotedblright{} of a continuous family of elliptic differential operators, or more generally of a \textquotedblleft relative elliptic complex\textquotedblright. For the initiated reader, let us say that the definition given here has the advantage over the classical definitions (see for example [2][17]) of requiring no deformation of the elliptic complex under consideration.

Sections 1 and 2 consist of reminders about the classical notion of a \textquotedblleft mixed structure\textquotedblright\SGAAuthorFootnote{1}{The informed reader is advised to skip them.}. We give, of course, only the minimum of definitions and results required to state and prove, in \S3, the central theorem we have in view, which says, a little more precisely, that the direct image, under the projection (assumed proper) onto the base (assumed locally compact), of a relative elliptic complex is a perfect complex. When the base $S$ is compact, the definition of the analytic index as an element of $K(S)$ (\textquotedblleft Grothendieck group\textquotedblright{} of complex vector bundles of finite rank on $S$) then follows from the globalization sorite of (II 2.3.2).

Throughout what follows, for a topological space $S$, we shall denote by $\underline{O}_S$ the sheaf of continuous real-valued functions on $S$.

\sgaStableTarget{sga6.e02.sec.1.occ3}\subsection*{1. The sorite of mixed structures}

In this section, we fix a topological space $S$.

\sgaStableTarget{sga6.e02.sec.1.1}\subsubsection*{1.1. Definition of mixed varieties}
Let $V$ be an open subset of $\underline{\underline{R}}^n$. For $r\in\underline{\underline{N}}$, we denote by $C^r(V)$ the Fréchet space of real-valued functions of class $C^r$ on $V$. We put
\[
  C^\infty(V)=\varprojlim_r C^r(V).
\]
Let $U$ be an open subset of $S$. For $r\in\underline{\underline{N}}\cup\{\infty\}$, we denote by ${}^r\underline{O}_{U\times V}$ the sheaf of rings on $U\times V$ defined by
%%%%%%%%%%  scan idx 213  |  folio 200  %%%%%%%%%%
\SGAPrintPage{200}%
\[
  \Gamma(U'\times V',{}^r\underline{O}_{U'\times V'})=C(U',C^r(V')),
\]
where $U'$ (resp. $V'$) is an open subset of $U$ (resp. $V$), and $C(U',C^r(V'))$ denotes the space of continuous maps from $U'$ to $C^r(V')$. There is a canonical injection
\[
  \pr_1^{-1}(\underline{O}_U)\longrightarrow{}^r\underline{O}_{U\times V},
\]
which makes $(U\times V,{}^r\underline{O}_{U\times V})$ a ringed space over $U$.

A ringed space $X$ over $S$ is called a \uline{mixed variety of class} $C^r$ ($r\in\underline{\underline{N}}\cup\{\infty\}$) when $X$ is locally isomorphic (as a ringed space over $S$) to a space of the form $(U\times V,{}^r\underline{O}_{U\times V})$. (When $S$ consists of one point, a mixed variety of class $C^r$ over $S$ is a variety of class $C^r$ in the usual sense).

Let $f:X\to S$ be a mixed variety of class $C^r$. For $s\in S$, the fiber $X_s=f^{-1}(s)$, equipped with the sheaf $i_s^{-1}(\underline{O}_X)/I_s$, where $i_s:X_s\to X$ is the injection and $I_s$ is the Ideal of sections that \textquotedblleft vanish\textquotedblright{} on $X_s$, is a variety of class $C^r$.

\sgaStableTarget{sga6.e02.sec.1.2}\subsubsection*{1.2. \textquotedblleft Relative\textquotedblright{} differential operators}
It is easy to make an infinitesimal study of mixed varieties on the model of Algebraic Geometry (EGA IV 16, SGA 3 VII). We shall limit ourselves to a brief review of the notions we shall need.

\sgaStableTarget{sga6.e02.para.1-2-1}1.2.1. Let $f:X\to S$ be a mixed variety of class $C^\infty$ (1.1). For every $n\in\underline{\underline{N}}$, we define the sheaf $P^n_{X/S}$ of \uline{principal parts of order} $n$ on $X$. In two different ways, corresponding to the two projections $\pr_1$ and $\pr_2$ from $X\times_S X$ onto $X$, it is an $\underline{O}_X$-Algebra augmented to $\underline{O}_X$, locally free of finite type as an $\underline{O}_X$-Module. When we do not specify the $\underline{O}_X$-Module structure considered on $P^n_{X/S}$, it will be understood to be the one coming from $\pr_1$. There are canonical isomorphisms
%%%%%%%%%%  scan idx 214  |  folio 201  %%%%%%%%%%
\SGAPrintPage{201}%
\[
  \begin{aligned}
  P^0_{X/S}&\simeq\underline{O}_X,\\
  P^1_{X/S}&\simeq\underline{O}_X\oplus\Omega^1_{X/S},
  \end{aligned}
\]
where $\Omega^1_{X/S}$ is the sheaf of relative differentials of $X$ over $S$, or \uline{relative cotangent sheaf}. We denote by $T_{X/S}$ the dual of $\Omega^1_{X/S}$, called the \uline{relative tangent sheaf}. The sheaf $T_{X/S}$ (resp. $\Omega^1_{X/S}$) is the sheaf of sections of a vector bundle of finite rank on $X$, called the \uline{relative} \uline{tangent bundle} (resp. \uline{cotangent bundle}), which is a mixed $S$-variety over $X$.

The $f^{-1}(\underline{O}_S)$-linear homomorphism (corresponding to the $\underline{O}_X$-Algebra structure defined by $\pr_2$)
\[
  d^n_{X/S}:\underline{O}_X\longrightarrow P^n_{X/S},
\]
which associates to a section of $\underline{O}_X$ its principal part of order $n$, is called the \uline{universal differential operator of order} $n$ on $X$ relative to $S$.

\sgaStableTarget{sga6.e02.para.1-2-2}1.2.2. More generally, for a locally free $\underline{O}_X$-Module $E$ of finite type, we can consider the sheaf $P^n_{X/S}(E)=P^n_{X/S}\otimes_{\underline{O}_X}E$ of principal parts of order $n$ of $E$ (the tensor product being defined by means of the second $\underline{O}_X$-Module structure on $P^n_{X/S}$), and we have a universal differential operator
\[
  d^n_{X/S}(E):E\longrightarrow P^n_{X/S}(E).
\]

Let $E,F$ be locally free $\underline{O}_X$-Modules of finite type. By definition, a differential operator from $E$ to $F$ of order $\le n$ is an $f^{-1}(\underline{O}_S)$-linear homomorphism from $E$ to $F$ that factors through $d^n_{X/S}(E)$ as an $\underline{O}_X$-linear homomorphism from $P^n_{X/S}(E)$ to $F$ (the factorization is then unique). In other words, the operator $d^n_{X/S}(E)$ gives a bijection
\sgaStableTarget{sga6.e02.eq.1-2-2-1}\begin{equation}\tag{1.2.2.1}
  \Homop_{\underline{O}_X}\bigl(P^n_{X/S}(E),F\bigr)
  \xrightarrow{\sim}
  \Diff^n(E,F),
\end{equation}
where $\Diff^n(E,F)$ denotes the set of differential operators from $E$
%%%%%%%%%%  scan idx 215  |  folio 202  %%%%%%%%%%
\SGAPrintPage{202}%
to $F$ of order $\le n$.

\sgaStableTarget{sga6.e02.para.1-2-3}1.2.3. We shall also have to consider differential operators with complex coefficients. Denote by $\underline{O}_S\otimes\underline{\underline{C}}$ (resp. $\underline{O}_X\otimes\underline{\underline{C}}$) the complexification of $\underline{O}_S$ (resp. $\underline{O}_X$). Let $E,F$ be locally free $(\underline{O}_X\otimes\underline{\underline{C}})$-Modules of finite type. By a differential operator from $E$ to $F$ of order $\le n$, we shall mean an $f^{-1}(\underline{O}_S\otimes\underline{\underline{C}})$-linear homomorphism from $E$ to $F$ that factors through $d^n_{X/S}(E)$ as an $(\underline{O}_X\otimes\underline{\underline{C}})$-linear homomorphism from $P^n_{X/S}(E)$ to $F$. Denoting the set of these operators by $\Diff^n_{\underline{\underline{C}}}(E,F)$, the operator $d^n_{X/S}(E)$ therefore defines a bijection
\sgaStableTarget{sga6.e02.eq.1-2-3-1}\begin{equation}\tag{1.2.3.1}
  \Homop_{\underline{O}_X\otimes\underline{\underline{C}}}\bigl(P^n_{X/S}(E),F\bigr)
  \xrightarrow{\sim}
  \Diff^n_{\underline{\underline{C}}}(E,F).
\end{equation}

\sgaStableTarget{sga6.e02.para.1-2-4}1.2.4. For $n\ge1$, there is a canonical exact sequence
\[
  0\longrightarrow S^n(\Omega^1_{X/S})
  \longrightarrow P^n_{X/S}
  \longrightarrow P^{n-1}_{X/S}
  \longrightarrow 0.\text{\SGAAuthorFootMark{1}}
\]
\SGAAuthorFootText{1}{$S^n(\ )$ denotes the $n$th symmetric power.}%
Hence, tensoring over $\underline{O}_X$\SGAAuthorFootnote{2}{For the second $\underline{O}_X$-Module structure of $P^n_{X/S}$ and $P^{n-1}_{X/S}$.} with a locally free $\underline{O}_X$-Module (resp. $(\underline{O}_X\otimes\underline{\underline{C}})$-Module) $E$ of finite type, we obtain an exact sequence of $\underline{O}_X$ (resp. $(\underline{O}_X\otimes\underline{\underline{C}})$)-Modules:
\sgaStableTarget{sga6.e02.eq.1-2-4-1}\begin{equation}\tag{1.2.4.1}
  0\longrightarrow S^n(\Omega^1_{X/S})\otimes_{\underline{O}_X}E
  \longrightarrow P^n_{X/S}(E)
  \longrightarrow P^{n-1}_{X/S}(E)
  \longrightarrow 0.
\end{equation}
Let $F$ be a locally free $\underline{O}_X$ (resp. $(\underline{O}_X\otimes\underline{\underline{C}})$)-Module of finite type. Applying the functor $\Homop_{\underline{O}_X}(.,F)$ (resp. $\Homop_{\underline{O}_X\otimes\underline{\underline{C}}}(.,F)$) to (1.2.4.1), we deduce, taking account of the identifications (1.2.2.1) and (1.2.3.1), an exact sequence
\sgaStableTarget{sga6.e02.eq.1-2-4-2}\begin{equation}\tag{1.2.4.2}
\begin{gathered}
  0\longrightarrow \Diff^{n-1}(E,F)
  \longrightarrow \Diff^n(E,F)
  \xrightarrow{\sigma_n}
  \Homop_{\underline{O}_X}\bigl(S^n(\Omega^1_{X/S})\otimes_{\underline{O}_X}E,F\bigr)\\
  (\text{resp. } 0\longrightarrow \Diff^{n-1}_{\underline{\underline{C}}}(E,F)
  \longrightarrow \Diff^n_{\underline{\underline{C}}}(E,F)
  \longrightarrow
  \Homop_{\underline{O}_X\otimes\underline{\underline{C}}}\bigl(S^n(\Omega^1_{X/S})\otimes_{\underline{O}_X}E,F\bigr)).
\end{gathered}
\end{equation}
%%%%%%%%%%  scan idx 216  |  folio 203  %%%%%%%%%%
\SGAPrintPage{203}%
(When $X$ is paracompact and $\underline{O}_X$ is soft (cf. II 2.3.1), we can complete the exact sequence (1.2.4.2) by adding a zero on the right). The homomorphism $\sigma_n$ is called the \uline{symbol}; it associates to a differential operator from $E$ to $F$ of order $\le n$ a \textquotedblleft homogeneous polynomial map of degree $n$ from $\Omega^1_{X/S}$ to $\underline{\Homop}_{\underline{O}_X}(E,F)$ (resp. $\underline{\Homop}_{\underline{O}_X\otimes\underline{\underline{C}}}(E,F)$)\textquotedblright, which we may also regard, if we wish, as a homomorphism $p^*(E)\longrightarrow p^*(F)$ (where $p$ is the projection onto $X$ from the relative cotangent bundle (1.2.1)), satisfying a certain homogeneity condition.

Let $E,F,G$ be locally free $\underline{O}_X$-Modules of finite type, and $u\in\Diff^m(E,F)$, $v\in\Diff^n(F,G)$, whence $vu\in\Diff^{m+n}(E,G)$. We easily check the relation
\sgaStableTarget{sga6.e02.eq.1-2-4-3}\begin{equation}\tag{1.2.4.3}
  \sigma_{m+n}(vu)=\sigma_n(v)\sigma_m(u).
\end{equation}
(The same relation holds for differential operators with complex coefficients). In particular, when the composite of two differential operators is zero, the composite of their symbols is also zero.

\sgaStableTarget{sga6.e02.para.1-2-5}1.2.5. Let $E$ be a complex, with bounded degrees, of differential operators with complex coefficients on $X$ relative to $S$ (1.2.3), with $d^i:E^i\longrightarrow E^{i+1}$ of order $\le n_i$. By (1.2.4.3), we have
\[
  \sigma_{n_{i+1}}(d^{i+1})\,\sigma_{n_i}(d^i)=0.
\]
We call $E$ an \uline{elliptic complex} of order $\underline{n}=(n_i)$ when the complex $(p^*(E^i),\sigma_{n_i}(d^i))$ (where $p$ is the projection onto $X$ from the cotangent bundle) is acyclic outside the zero section.

When $S$ consists of one point, the preceding notion of an elliptic complex on $X$ coincides with the usual notion, studied in detail in [4] or [16].

\sgaStableTarget{sga6.e02.para.1-2-6}1.2.6. Let $s\in S$, and let $i_s$ be the inclusion of the fiber $X_s$ into $X$ (considered as a monomorphism of ringed spaces) cf.~(1.1). For $n\in\underline{\underline{N}}$, there is a canonical isomorphism

